% ============================================================
%  GUIA DE PREPARACION PEP 1 -- CALCULO I
%  Usach Premium
% ============================================================

\documentclass[letterpaper, 12pt]{article}

% ============================================================
% PAQUETES
% ============================================================
\usepackage[utf8]{inputenc}
\usepackage[spanish]{babel}
\usepackage{amsmath, amssymb, amsfonts}
\usepackage{geometry}
\usepackage{fancyhdr}
\usepackage{enumitem}
\usepackage{graphicx}
\usepackage{hyperref}
\usepackage{multicol}
\usepackage{needspace}
\usepackage{parskip}
\usepackage{xcolor}
\usepackage{tcolorbox}
\usepackage{eso-pic}
\usepackage{tikz}
\usetikzlibrary{patterns}
\tcbuselibrary{skins, breakable}

% ============================================================
% GEOMETRIA
% ============================================================
\geometry{letterpaper, margin=1.5cm, top=3cm, bottom=2.5cm, headheight=2cm}

% ============================================================
% COLORES
% ============================================================
\definecolor{celesteSuave}{HTML}{F0F7FF}
\definecolor{azulFuerte}{HTML}{1A5276}
\definecolor{turquesaUsach}{HTML}{33CCCC}
\definecolor{enlace}{HTML}{1A5276}
\definecolor{grisTexto}{HTML}{555555}

% ============================================================
% COMANDOS
% ============================================================
\newcommand{\R}{\mathbb{R}}
\newcommand{\Dom}{\operatorname{Dom}}
\newcommand{\Rec}{\operatorname{Rec}}
\newcommand{\Cod}{\operatorname{Cod}}

% ============================================================
% RUTAS DE IMAGENES
% ============================================================
\newcommand{\logoup}{../../../../../template/images/logo_up.png}
\newcommand{\logousach}{../../../../../template/images/logo_usach.png}
\newcommand{\logofondo}{../../../../../template/images/logo_up.png}

% ============================================================
% INFORMACION DE LA GUIA
% ============================================================
\newcommand{\institucion}{Usach Premium}
\newcommand{\curso}{C\'alculo I}
\newcommand{\guiasnumero}{Gu\'ia de Preparaci\'on PEP 1}
\newcommand{\semestre}{1er. Semestre de 2026}
\newcommand{\preparadopor}{Jaime Riquelme -- Usach Premium.}
\newcommand{\webinstitucion}{www.usachpremium.cl}
\newcommand{\instagraminstitucion}{https://www.instagram.com/usach.premium}
\newcommand{\transparencia}{0.05}
\newcommand{\fraseportada}{``Por y para estudiantes''}

\newcommand{\listacontenidos}{
    \item[\color{azulFuerte}$\Rightarrow$] N\'umeros reales, intervalos e inecuaciones
    \item[\color{azulFuerte}$\Rightarrow$] Valor absoluto y tablas de signos
    \item[\color{azulFuerte}$\Rightarrow$] Dominio, recorrido y composici\'on de funciones
    \item[\color{azulFuerte}$\Rightarrow$] Inyectividad, sobreyectividad y biyectividad
    \item[\color{azulFuerte}$\Rightarrow$] Funci\'on inversa y problemas tipo PEP
}

% ============================================================
% HYPERREF
% ============================================================
\hypersetup{
    colorlinks=true,
    linkcolor=enlace,
    urlcolor=enlace,
    citecolor=enlace,
    pdfborder={0 0 0},
}

% ============================================================
% ENCABEZADO Y PIE
% ============================================================
\pagestyle{fancy}
\fancyhf{}
\fancyhead[L]{%
    \includegraphics[height=1.2cm]{\logoup}%
}
\fancyhead[R]{%
    \begin{minipage}[b]{0.42\textwidth}
        \raggedleft
        \fontsize{9pt}{11pt}\selectfont
        \textbf{\color{azulFuerte}\institucion} \\
        \curso\ -- \guiasnumero \\
        \textit{\color{gray}@usach.premium}
    \end{minipage}\hspace{0.1cm}%
    \includegraphics[height=1.2cm]{\logousach}%
}
\fancyfoot[L]{\fontsize{9pt}{11pt}\selectfont \textit{\fraseportada}}
\fancyfoot[R]{\fontsize{9pt}{11pt}\selectfont P\'agina \thepage}
\renewcommand{\headrulewidth}{0.4pt}
\renewcommand{\footrulewidth}{0.4pt}

\fancypagestyle{plain}{
    \fancyhf{}
    \renewcommand{\headrulewidth}{0pt}
    \renewcommand{\footrulewidth}{0pt}
}

% ============================================================
% MARCA DE AGUA
% ============================================================
\newcommand{\MarcaAgua}{
    \AddToShipoutPicture{
        \AtPageCenter{
            \begin{tikzpicture}[remember picture, overlay]
                \node[opacity=\transparencia] at (0,0)
                    {\includegraphics[width=0.7\textwidth]{\logofondo}};
            \end{tikzpicture}%
        }
    }
}

% ============================================================
% CAJAS
% ============================================================
\newenvironment{kitpropiedades}[1]{%
    \begin{tcolorbox}[colback=celesteSuave, colframe=azulFuerte,
        title=\textbf{#1}, fonttitle=\bfseries, breakable]
}{%
    \end{tcolorbox}
}

\newtcolorbox{pepbox}[1]{
    colback=white,
    colframe=azulFuerte,
    fonttitle=\bfseries,
    coltitle=white,
    title=#1,
    arc=5pt,
    boxrule=1.1pt,
    enhanced,
    before={\par\Needspace{0.82\textheight}},
    drop shadow={black!5}
}

\newenvironment{solbox}[1]{%
    \begin{tcolorbox}[colback=celesteSuave, colframe=azulFuerte,
        title=\textbf{#1}, fonttitle=\bfseries]
}{%
    \end{tcolorbox}
}

\newenvironment{soluciones}{%
    \newpage
    \begin{center}
        \begin{tcolorbox}[colback=azulFuerte, colframe=azulFuerte,
            colupper=white, width=0.8\textwidth, center, arc=10pt]
            \centering\Large\textbf{SOLUCIONARIO}
        \end{tcolorbox}
    \end{center}
    \MarcaAgua
    \small
}{}

% ============================================================
\begin{document}
\thispagestyle{plain}

% ============================================================
% PORTADA
% ============================================================
\AddToShipoutPicture*{%
    \put(54, 700){\includegraphics[width=2cm]{\logoup}}
    \put(420, 706){\includegraphics[width=5cm]{\logousach}}
}

\vspace*{0.5cm}

\begin{center}
    \color{azulFuerte}
    {\huge \textbf{GU\'IA DE EJERCICIOS}} \\[0.5cm]
    {\Huge \textbf{\curso}} \\[0.3cm]
    {\Large \guiasnumero} \\[1.5cm]

    \begin{tcolorbox}[colback=celesteSuave, colframe=azulFuerte, arc=10pt,
        width=0.85\textwidth, center, boxrule=1.5pt]
        \centering
        \vspace{0.2cm}
        {\Large \textbf{Contenidos PEP 1}}
        \vspace{0.3cm}
        \color{black}
        \begin{itemize}[leftmargin=2.6cm]
            \listacontenidos
        \end{itemize}
        \vspace{0.2cm}
    \end{tcolorbox}
\end{center}

\vspace{1.5cm}

\begin{center}
    {\Large \textbf{\semestre}} \\[0.8cm]
    {\large \textbf{\preparadopor}}
\end{center}

\vspace{2cm}

\begin{center}
    {\color{azulFuerte}\hrule height 0.5pt}
    \vspace{0.3cm}
    \begin{minipage}{0.45\textwidth}
        \centering
        \textbf{Instagram:} \\
        \small\href{\instagraminstitucion}{@usach.premium}
    \end{minipage}
    \begin{minipage}{0.45\textwidth}
        \centering
        \textbf{Web:} \\
        \small\href{http://\webinstitucion}{\webinstitucion}
    \end{minipage}
\end{center}

\pagenumbering{arabic}
\newpage
\MarcaAgua

% ============================================================
% RESUMEN
% ============================================================
\begin{kitpropiedades}{LO QUE DEBES TENER A MANO}
\begin{multicols}{2}
\textbf{1. Inecuaciones.}
\begin{itemize}[leftmargin=*]
    \item Factoriza antes de hacer tabla de signos.
    \item Los ceros del denominador nunca se incluyen.
    \item En ra\'ices pares: el radicando debe ser $\geq 0$.
    \item Si elevas al cuadrado una desigualdad, verifica que ambos lados sean no negativos.
\end{itemize}

\columnbreak

\textbf{2. Funciones.}
\begin{itemize}[leftmargin=*]
    \item $\Dom(f\circ g)=\{x\in\Dom(g):g(x)\in\Dom(f)\}$.
    \item Para que exista $f^{-1}$, la funci\'on debe ser biyectiva.
    \item Para hacer biyectiva una par\'abola, restringe el dominio a un lado del v\'ertice.
    \item Si $\Cod(f)=\Rec(f)$ y $f$ es inyectiva, entonces $f$ es biyectiva.
\end{itemize}
\end{multicols}
\end{kitpropiedades}

% ============================================================
% EJERCICIOS
% ============================================================
\section*{I. N\'umeros reales e inecuaciones}

\begin{enumerate}[label=\color{azulFuerte}\bfseries 1.\arabic*., leftmargin=*, itemsep=1em]
    \item Resuelva la inecuaci\'on
    \[
        \frac{x-4}{x^{2}-9}\geq 0.
    \]

    \item Determine el conjunto soluci\'on de
    \[
        \frac{|x-1|-3}{x+2}<0.
    \]

    \item Resuelva en $\R$:
    \[
        \frac{x+2}{x-1}\leq \frac{x-3}{x+4}.
    \]

    \item El largo de un terreno rectangular supera al ancho en $7$ metros. Determine las dimensiones posibles si el \'area debe ser a lo m\'as $260\,m^2$.

    \item Resuelva la inecuaci\'on con ra\'iz:
    \[
        \sqrt{x+3}\leq x-1.
    \]

    \item Determine el conjunto soluci\'on de
    \[
        \sqrt{x^2-4x}>x-1.
    \]
\end{enumerate}

\section*{II. Dominio, recorrido y composici\'on}

\begin{enumerate}[label=\color{azulFuerte}\bfseries 2.\arabic*., leftmargin=*, itemsep=1em]
    \item Halle el dominio natural de
    \[
        h(x)=\sqrt{\frac{x+4}{2-x}}.
    \]

    \item Considere $f(x)=\sqrt{5-x}$ y $g(x)=\sqrt{x^2-1}$. Determine $\Dom(f\circ g)$.

    \item Considere
    \[
        f(x)=\sqrt{\frac{x-1}{x+3}},
        \qquad
        g(x)=2x-5.
    \]
    Determine $\Dom(f\circ g)$, halle expl\'icitamente $(f\circ g)(x)$, calcule $(f\circ g)(4)$ y determine si existe una preimagen de $1$.

    \item Sea
    \[
        f(x)=1+\sqrt{\frac{x+1}{x-2}}.
    \]
    Determine $\Dom(f)$ y la preimagen de $3$, si existe.

    \item Considere
    \[
        p(x)=\sqrt{6-\sqrt{x+2}}.
    \]
    Determine $\Dom(p)$, indique si $p$ es creciente o decreciente en su dominio, y halle $\Rec(p)$.

    \item Sea
    \[
        f(x)=\sqrt{\frac{4}{\sqrt{x-2}}-1},
        \qquad
        g(x)=x+3.
    \]
    Halle $\Dom(f)$, $\Rec(f)$ y $\Dom(f\circ g)$, indicando adem\'as la regla de asignaci\'on de $f\circ g$.
\end{enumerate}

\section*{III. Inyectividad, biyectividad e inversa}

\begin{enumerate}[label=\color{azulFuerte}\bfseries 3.\arabic*., leftmargin=*, itemsep=1em]
    \item Sea $f:A\to B$ definida por
    \[
        f(x)=(x+1)^2-2.
    \]
    Determine conjuntos $A$ y $B$ para que $f$ sea biyectiva y $2\in A$. Luego encuentre $f^{-1}$.

    \item Sea $k\in\R$ tal que
    \[
        f:[-2,1]\to[k,27],
        \qquad
        f(x)=(x-3)^2+2
    \]
    est\'a bien definida. Demuestre que $f$ es inyectiva, determine $k$ para que sea sobreyectiva y encuentre $f^{-1}$ para ese valor de $k$.

    \item Considere
    \[
        q(x)=2-\sqrt{9-(x-4)^2}.
    \]
    Determine $\Dom(q)$. Luego escoja un subconjunto $A_1\subseteq \Dom(q)$ donde $q$ sea inyectiva y encuentre una funci\'on inversa indicando dominio y recorrido.

    \item Considere
    \[
        F:]-2,2[\to\R,
        \qquad
        F(x)=\frac{x}{2-|x|}.
    \]
    Pruebe que $F$ es biyectiva y determine $F^{-1}$.

    \item Sea
    \[
        r:\R-\{-2\}\to\R-\{3\},
        \qquad
        r(x)=\frac{3x-1}{x+2}.
    \]
    Justifique que $r$ es biyectiva y encuentre $r^{-1}$.
\end{enumerate}

\section*{IV. Par\'ametros e intersecciones}

\begin{enumerate}[label=\color{azulFuerte}\bfseries 4.\arabic*., leftmargin=*, itemsep=1em]
    \item Determine todos los valores de $k\in\R$ para que la par\'abola
    \[
        y=kx^2+4x+1
    \]
    y la recta
    \[
        y=(k+2)x-3
    \]
    se intersecten en dos puntos diferentes.

    \item Considere las par\'abolas
    \[
        C_1:\ y=(k+1)x^2-3x+k,
        \qquad
        C_2:\ y=2x^2+(k-2)x+1.
    \]
    Determine los valores de $k\in\R$ para que $C_1$ y $C_2$ se intersecten en dos puntos diferentes.
\end{enumerate}

\newpage
\section*{V. Simulacros tipo PEP}

\begin{pepbox}{Simulacro A}
\begin{enumerate}[label=\textbf{\arabic*.}, leftmargin=*, itemsep=1em]
    \item \textbf{[20 puntos]} Considere
    \[
        f(x)=1+\sqrt{\frac{x-3}{x+2}}.
    \]
    \begin{enumerate}[label=\textbf{\alph*)}]
        \item Determine $\Dom(f)$.
        \item Determine $\Dom(f\circ f)$.
    \end{enumerate}

    \item \textbf{[20 puntos]} Considere $g:A\to B$ definida por
    \[
        g(x)=-2+\sqrt{-x^2+8x-12}.
    \]
    \begin{enumerate}[label=\textbf{\alph*)}]
        \item Determine $A=\Dom(g)$.
        \item Determine un conjunto $A_1\subseteq A$ para que $g$ sea inyectiva y demu\'estrelo.
        \item Determine $B$ para que $g:A_1\to B$ sea biyectiva y encuentre $g^{-1}$.
    \end{enumerate}

    \item \textbf{[20 puntos]} Determine los valores de $k\in\R$ para que la par\'abola
    \[
        y=(k-1)x^2+2x+k
    \]
    y la recta
    \[
        y=kx-4
    \]
    se intersecten en dos puntos distintos.
\end{enumerate}
\end{pepbox}

\begin{pepbox}{Simulacro B}
\begin{enumerate}[label=\textbf{\arabic*.}, leftmargin=*, itemsep=1em]
    \item \textbf{[20 puntos]} Resuelva la siguiente inecuaci\'on:
    \[
        \frac{3-|x+1|}{x^2-4}<0.
    \]

    \item \textbf{[20 puntos]} Considere
    \[
        f(x)=\sqrt{\frac{5}{\sqrt{x+1}}-1},
        \qquad
        g(x)=x-4.
    \]
    \begin{enumerate}[label=\textbf{\alph*)}]
        \item Halle $\Dom(f)$.
        \item Halle $\Rec(f)$.
        \item Determine $f\circ g$ indicando expresamente su dominio y regla de asignaci\'on.
    \end{enumerate}

    \item \textbf{[20 puntos]} Sea $k\in\R$ tal que
    \[
        f:[-2,0]\to[k,12],
        \qquad
        f(x)=(x-1)^2+3
    \]
    est\'a bien definida.
    \begin{enumerate}[label=\textbf{\alph*)}]
        \item Demuestre que $f$ es inyectiva.
        \item Determine $k$ para que $f$ sea sobreyectiva.
        \item Para ese valor de $k$, determine $f^{-1}$ indicando dominio, codominio y regla.
    \end{enumerate}
\end{enumerate}
\end{pepbox}

% ============================================================
% SOLUCIONARIO
% ============================================================
\begin{soluciones}

\begin{solbox}{I. Soluciones -- Inecuaciones}
\begin{enumerate}[label=\color{azulFuerte}\bfseries 1.\arabic*., leftmargin=*, itemsep=0.55em]
    \item $x\in(-3,3)\cup[4,+\infty)$.
    \item $x\in(-\infty,-2)\cup(-2,4)$.
    \item $x\in(-\infty,-4)\cup\left[-\dfrac12,1\right)$.
    \item Si $x$ es el ancho, $x(x+7)\leq260$, luego $x\in[0,13]$. Ancho entre $0$ y $13$ m; largo entre $7$ y $20$ m.
    \item $x\in\left[\dfrac{3+\sqrt{17}}{2},+\infty\right)$.
    \item $x\in(-\infty,0]$.
\end{enumerate}
\end{solbox}

\begin{solbox}{II. Soluciones -- Dominio, recorrido y composici\'on}
\begin{enumerate}[label=\color{azulFuerte}\bfseries 2.\arabic*., leftmargin=*, itemsep=0.55em]
    \item $\Dom(h)=[-4,2)$.
    \item $\Dom(f\circ g)=[-\sqrt{26},-1]\cup[1,\sqrt{26}]$.
    \item $\Dom(f\circ g)=(-\infty,1)\cup[3,+\infty)$,
    \[
        (f\circ g)(x)=\sqrt{\frac{x-3}{x-1}}.
    \]
    Adem\'as $(f\circ g)(4)=\sqrt{\frac13}$ y no existe preimagen de $1$.
    \item $\Dom(f)=(-\infty,-1]\cup(2,+\infty)$; la preimagen de $3$ es $x=3$.
    \item $\Dom(p)=[-2,34]$; $p$ es decreciente; $\Rec(p)=[0,\sqrt6]$.
    \item $\Dom(f)=(2,18]$, $\Rec(f)=[0,+\infty)$,
    \[
        \Dom(f\circ g)=(-1,15],
        \qquad
        (f\circ g)(x)=\sqrt{\frac{4}{\sqrt{x+1}}-1}.
    \]
\end{enumerate}
\end{solbox}

\begin{solbox}{III. Soluciones -- Biyectividad e inversa}
\begin{enumerate}[label=\color{azulFuerte}\bfseries 3.\arabic*., leftmargin=*, itemsep=0.65em]
    \item Una elecci\'on natural es
    \[
        A=[-1,+\infty),\qquad B=[-2,+\infty).
    \]
    Entonces
    \[
        f^{-1}(y)=-1+\sqrt{y+2},
        \qquad
        f^{-1}:[-2,+\infty)\to[-1,+\infty).
    \]

    \item En $[-2,1]$, la funci\'on $(x-3)^2+2$ es estrictamente decreciente, luego es inyectiva. Su recorrido es $[6,27]$, por lo tanto $k=6$. La inversa es
    \[
        f^{-1}(y)=3-\sqrt{y-2},
        \qquad
        f^{-1}:[6,27]\to[-2,1].
    \]

    \item $\Dom(q)=[1,7]$. Tomando $A_1=[4,7]$, se tiene $\Rec(q)=[-1,2]$ y
    \[
        q^{-1}(y)=4+\sqrt{9-(y-2)^2},
        \qquad
        q^{-1}:[-1,2]\to[4,7].
    \]

    \item La funci\'on es estrictamente creciente en cada tramo $]-2,0[$ y $[0,2[$, con recorrido $(-\infty,0)$ y $[0,+\infty)$ respectivamente. Por tanto es biyectiva. Su inversa es
    \[
        F^{-1}(y)=\frac{2y}{1+|y|},
        \qquad
        F^{-1}:\R\to]-2,2[.
    \]

    \item Si $y=\dfrac{3x-1}{x+2}$, entonces
    \[
        x=\frac{1+2y}{3-y}.
    \]
    As\'i,
    \[
        r^{-1}(y)=\frac{1+2y}{3-y},
        \qquad
        r^{-1}:\R-\{3\}\to\R-\{-2\}.
    \]
\end{enumerate}
\end{solbox}

\begin{solbox}{IV. Soluciones -- Par\'ametros}
\begin{enumerate}[label=\color{azulFuerte}\bfseries 4.\arabic*., leftmargin=*, itemsep=0.65em]
    \item Al igualar:
    \[
        kx^2+(2-k)x+4=0.
    \]
    Debe cumplirse $k\neq0$ y
    \[
        \Delta=(2-k)^2-16k=k^2-20k+4>0.
    \]
    Por lo tanto,
    \[
        k\in(-\infty,0)\cup(0,10-4\sqrt6)\cup(10+4\sqrt6,+\infty).
    \]

    \item Al igualar:
    \[
        (k-1)x^2-(k+1)x+(k-1)=0.
    \]
    Debe cumplirse $k\neq1$ y
    \[
        \Delta=(k+1)^2-4(k-1)^2=-3k^2+10k-3>0.
    \]
    Luego,
    \[
        k\in\left(\frac13,1\right)\cup(1,3).
    \]
\end{enumerate}
\end{solbox}

\begin{solbox}{V. Soluciones -- Simulacro A}
\begin{enumerate}[label=\color{azulFuerte}\bfseries A.\arabic*., leftmargin=*, itemsep=0.8em]
    \item
    \[
        \Dom(f)=(-\infty,-2)\cup[3,+\infty).
    \]
    Como $f(x)\geq1$, para que $f(x)\in\Dom(f)$ basta exigir $f(x)\geq3$. Entonces
    \[
        1+\sqrt{\frac{x-3}{x+2}}\geq3
        \Longleftrightarrow
        \frac{x-3}{x+2}\geq4
        \Longleftrightarrow
        \frac{3x+11}{x+2}\leq0.
    \]
    Por lo tanto,
    \[
        \Dom(f\circ f)=\left[-\frac{11}{3},-2\right).
    \]

    \item
    \[
        -x^2+8x-12\geq0
        \Longleftrightarrow
        (x-2)(x-6)\leq0,
    \]
    luego $A=[2,6]$. Se puede escoger $A_1=[4,6]$. En ese intervalo $g$ es decreciente, y tambi\'en se puede demostrar por definici\'on: si $a,b\in[4,6]$ y $g(a)=g(b)$, entonces
    \[
        -a^2+8a-12=-b^2+8b-12
        \Longleftrightarrow
        (a-b)(a+b-8)=0.
    \]
    Como $a,b\in[4,6]$, la igualdad $a+b=8$ solo puede producir repetici\'on sim\'etrica respecto de $4$; en $[4,6]$ esto fuerza $a=b$ para mantener la misma rama. As\'i, $g$ es inyectiva en $[4,6]$.

    En $[4,6]$, $\Rec(g)=[-2,0]$. Por tanto $B=[-2,0]$ y
    \[
        g^{-1}(y)=4+\sqrt{4-(y+2)^2},
        \qquad
        g^{-1}:[-2,0]\to[4,6].
    \]

    \item Al igualar:
    \[
        (k-1)x^2+(2-k)x+(k+4)=0.
    \]
    Debe cumplirse $k\neq1$ y
    \[
        \Delta=(2-k)^2-4(k-1)(k+4)>0
        \Longleftrightarrow
        3k^2+16k-20<0.
    \]
    Luego,
    \[
        k\in\left(\frac{-8-2\sqrt{31}}{3},1\right)\cup
        \left(1,\frac{-8+2\sqrt{31}}{3}\right).
    \]
\end{enumerate}
\end{solbox}

\begin{solbox}{V. Soluciones -- Simulacro B}
\begin{enumerate}[label=\color{azulFuerte}\bfseries B.\arabic*., leftmargin=*, itemsep=0.8em]
    \item Los puntos cr\'iticos son $x=-4$, $x=-2$ y $x=2$; estos dos \'ultimos anulan el denominador. La tabla de signos entrega
    \[
        x\in(-\infty,-4)\cup(-2,2)\cup(2,+\infty).
    \]

    \item
    \[
        \Dom(f)=(-1,24],
        \qquad
        \Rec(f)=[0,+\infty).
    \]
    Adem\'as,
    \[
        \Dom(f\circ g)=\{x\in\R:x-4\in(-1,24]\}=(3,28],
    \]
    y
    \[
        (f\circ g)(x)=\sqrt{\frac{5}{\sqrt{x-3}}-1}.
    \]

    \item En $[-2,0]$, la funci\'on $(x-1)^2+3$ es estrictamente decreciente, luego es inyectiva. Su recorrido es
    \[
        \Rec(f)=[4,12],
    \]
    por lo que $k=4$. Para ese valor,
    \[
        f^{-1}(y)=1-\sqrt{y-3},
        \qquad
        f^{-1}:[4,12]\to[-2,0].
    \]
\end{enumerate}
\end{solbox}

\vspace{0.8cm}
\begin{center}
    \small\textbf{Usach Premium} -- Nos vemos en la ayudant\'ia. \\
    \textit{\fraseportada}
\end{center}

\end{soluciones}

\end{document}
